%%%% RESUMO

\begin{resumoutfpr}
Este trabalho desenvolve e avalia um protótipo acadêmico que integra um processador ARM e um arranjo de portas programáveis em campo (FPGA) para processar dados sintéticos de mercado e executar uma regra de decisão gerada por síntese de alto nível. Um gerador transmite pelo Protocolo de Controle de Transmissão (TCP) mensagens sintéticas codificadas segundo \textit{FIX Adapted for STreaming} (FAST), e um receptor C++ as decodifica e converte em quadros de 256 bits. No FPGA da placa DE10-Nano, filas circulares transportam os eventos, um livro de ofertas agregado mantém oito níveis por lado e uma função escrita em MATLAB, convertida para VHDL pelo HDL Coder, produz os códigos NOOP, BUY ou SELL. A ponte, o livro e os invólucros permanecem em VHDL manual. A validação compreende testes C++, bancos de teste VHDL e um ensaio funcional na placa. Em uma campanha registrada com um milhão de eventos, o núcleo FPGA apresentou latência de três ciclos (60~ns a 50~MHz), com variação desprezível na resolução de 20~ns do contador, enquanto o núcleo C++ de referência apresentou média de 315~ns. A latência média de ida e volta pela entrada e saída mapeada em memória foi 10.228~ns, com percentil 99 de 10.310~ns, e a vazão em fluxo foi 102.381 mensagens por segundo. Os valores caracterizam uma execução específica e não incluem TCP ou decodificação FAST. O protótipo não se conecta a bolsa, não envia ordens, não implementa controle de risco e não constitui um sistema operacional de negociação de alta frequência.
\end{resumoutfpr}
